\section{Introduction}
\label{sec:introduction}

An 8-bit delivery frame is the end of a lossy chain. Scene radiance spanning
perhaps twenty stops was metered, exposed, graded through a tone curve chosen by
a colourist, clipped at both ends and quantised to 256 levels per channel.
Single-image inverse tone mapping asks for the radiance back. The task is
attractive because the input is abundant and the output is what every high
dynamic range (HDR) pipeline needs, and it is hard for a reason that is easy to
state and easy to forget: the forward map is not injective, and no amount of
model does anything about that.

The literature has largely treated the problem as one of architecture and
training. We treat it as one of information, and the distinction turns out to be
load-bearing. This paper reports a compact model, the conditions under which it
helps and hurts, a measurement of the ceiling on adapting it per frame, and a
component that reaches a particular part of that ceiling by asking a different
question of the same input.

\paragraph{Position of this work.}
Our model is deliberately small (\num{1196197} parameters) and is a residual on
an analytic inverse tone map rather than a direct predictor. That choice is what
makes the paper's central question askable at every point in training and
evaluation: \emph{does the network beat doing nothing?} A direct predictor has
no ``nothing'' to be compared against, and in practice a great deal of what a
learned inverse tone mapper appears to contribute is recoverable analytically.
Every number in this paper is a difference against that analytic baseline on
identical inputs.

\paragraph{Contributions.}
\begin{enumerate}
\item \textbf{A conditional characterisation of when learned inverse tone
  mapping helps.} On degraded input the model gains \SI{1.43}{\decibel}
  PU21-PSNR and \num{0.44}~JOD over the analytic baseline, on 81\% of held-out
  frames. On clean input it loses \SI{3.0}{\decibel} PU21-PSNR for
  $-0.046$~JOD. Reporting either condition alone misrepresents the method
  (\cref{sec:results}).

\item \textbf{A characterised disagreement between the two standard HDR
  metrics.} On the same frames and the same models, PU21-PSNR reports a
  substantial fidelity penalty where CVVDP reports a perceptually negligible
  difference, and the two agree decisively when the difference is large. The
  two are not commensurable as numbers, which is the point: what differs is the
  practical significance each assigns to a small reconstruction difference on
  well-graded input. We localise that divergence, show it is only partly carried
  by the error tail, and show it reappearing inside a checkpoint-selection rule
  that touches neither metric
  (\cref{sec:results,sec:failure,sec:methodology}).

\item \textbf{A localisation of the failure to a specific mechanism.} The clean
  regression is not invented highlights, which was our own first two hypotheses.
  It is the shadow arm of the residual gate firing on scenes that need no
  reconstruction, and disabling that arm inverts the regression
  (\cref{sec:failure}).

\item \textbf{A measured bound on per-frame adaptation.} An oracle per-frame
  residual scale is worth \SI{+5.84}{\decibel} on clean and
  \SI{+0.29}{\decibel} on degraded input simultaneously, and is poorly
  identifiable from the input we have: a linear readout of the features the gate
  sees explains \num{3}\% of its variance, and \num{79}\% of the variance lies
  within a condition rather than between, which caps \emph{any}
  condition-detection architecture at $R^2 = \num{0.213}$ regardless of how it
  is built. Capacity, corpus and objective were each varied by a factor of four
  to six without moving it. We are careful to claim an empirical identifiability
  result and not an information-theoretic one; \cref{sec:bound} separates what
  the measurements establish from what they suggest.

\item \textbf{A component that reaches the reachable part of that bound.} Posed
  as a binary decision on the detectable axis rather than as a continuous scale,
  the problem is tractable. A \num{21121}-parameter gate on the shadow arm
  delivers $+3.41 \pm 0.32$~dB and $+0.135 \pm 0.023$~JOD on clean input
  for $0.30 \pm 0.15$~dB and $0.091 \pm 0.041$~JOD on degraded input
  across three seeds, and is the only configuration we scored that is positive
  on all four measures (\cref{sec:gate}).

\item \textbf{Three defects in our own evaluation harness, reported rather than
  quietly repaired.} Checkpoint selection on the maximum of a noisy series, an
  evaluation that silently read the alphabetical front of its split, and a
  viewer that presented every frame upside down beneath tests that compared only
  float buffers. Each corrupted a result we believed, and each is available to
  any comparable pipeline (\cref{sec:methodology}).
\end{enumerate}

\paragraph{Scope.}
This paper reports a direct SDR-pixel to scene-linear-radiance image model. It
does not report the DRE-transformer and cross-attention architecture described
in an earlier manuscript from this group: that design was never trained, and the
results attributed to it were aspirational. Both are withdrawn.
\Cref{sec:withdrawal} states what was withdrawn and why, so that a reader who
has seen the earlier document can reconcile the two. What follows is the system
that exists, measured with the metrics that exist.
